\documentclass[a4paper,13pt,twoside]{report}
\usepackage{scrextend}
\changefontsizes{13pt}
\usepackage[utf8]{vietnam}
\usepackage[top=2cm, bottom=2cm, left=3.5cm, right=2.5cm]{geometry}

\usepackage{graphicx} % Cho phép chèn hỉnh ảnh
\usepackage{fancybox} % Tạo khung box
\usepackage{indentfirst} % Thụt đầu dòng ở dòng đầu tiên trong đoạn
\usepackage{amsthm} % Cho phép thêm các môi trường định nghĩa
\usepackage{latexsym} % Các kí hiệu toán học
\usepackage{amsmath} % Hỗ trợ một số biểu thức toán học
\usepackage{amssymb} % Bổ sung thêm kí hiệu về toán học
\usepackage{amsbsy} % Hỗ trợ các kí hiệu in đậm
\usepackage{times} % Chọn font Time New Romans
\usepackage{array} % Tạo bảng array
\usepackage{longtable} % Bảng dài tự động sang trang (dùng cho danh mục viết tắt)
\usepackage{enumitem} % Cho phép thay đổi kí hiệu của list
\usepackage{subfiles} % Chèn các file nhỏ, giúp chia các chapter ra nhiều file hơn
\usepackage{titlesec} % Giúp chỉnh sửa các tiêu đề, đề mục như chương, phần,..
\usepackage{titletoc}
\usepackage{chngcntr} % Dùng để thiết lập lại cách đánh số caption,..
\usepackage{pdflscape} % Đưa các bảng có kích thước đặt theo chiều ngang giấy
\usepackage{rotating} % sidewaysfigure: xoay ngang sơ đồ ERD lớn để phóng to
\usepackage{afterpage}
\usepackage[ruled,vlined]{algorithm2e}  % Hỗ trợ viết các giải thuật
\usepackage{capt-of} % Cho phép sử dụng caption lớn đối với landscape page
\usepackage{multirow} % Merge cells
\usepackage{fancyhdr} % Cho phép tùy biến header và footer
% \usepackage[natbib,backend=biber,style=ieee]{biblatex} % Giúp chèn tài liệu tham khảo
\usepackage{appendix}

\usepackage[font=small,labelfont=bf]{caption}

\usepackage{listings}
\usepackage{float}
\usepackage{subcaption}
\usepackage{xurl}
\usepackage{outlines}

% Giảm khoảng trắng quanh hình/bảng (tránh khoảng trống lớn khi ảnh không vừa cuối trang)
% raggedbottom: không giãn khoảng cách dọc để căn đáy trang (twoside mặc định flushbottom gây khoảng trống lớn quanh hình/bảng)
\raggedbottom
\setlength{\intextsep}{6pt plus 2pt minus 2pt}
\setlength{\textfloatsep}{8pt plus 2pt minus 2pt}
\setlength{\floatsep}{6pt plus 2pt minus 2pt}
\setlength{\abovecaptionskip}{4pt}
\setlength{\belowcaptionskip}{2pt}
% Hạn chế dòng mồ côi (widow/orphan) đầu/cuối trang
\widowpenalty=10000
\clubpenalty=10000

\usepackage{setspace}
\usepackage{parskip}

% package content table
\usepackage{tocbasic}


% ===================================================

 




\renewcommand\appendixname{PHỤ LỤC}
\renewcommand\appendixpagename{PHỤ LỤC}
\renewcommand\appendixtocname{PHỤ LỤC}


 % chèn file chứa danh mục tài liệu tham khảo vào 

% Tên cột DB trong bảng schema — monospace, ngắt dòng sát _ không thêm trắng
\makeatletter
\newcommand{\dbid}[1]{{\ttfamily\@db@split#1_\@nil}}
\def\@db@split#1_#2\@nil{%
  #1%
  \ifx\@nil#2\@nil\else
    \textunderscore\discretionary{}{}{}\@db@split#2\@nil
  \fi}
\makeatother

% Cột bảng ragged-right, chặn ngắt giữa từ tiếng Anh (pass-word) nhưng vẫn cho
% phép ngắt ở dấu _ của \dbid (tránh tên cột dài chờm sang cột bên)
\newcolumntype{R}[1]{>{\raggedright\arraybackslash\hyphenpenalty=10000}p{#1}}

% Cho phép LaTeX nới rộng khoảng cách để tránh tràn margin với \texttt{} dài
\setlength{\emergencystretch}{3em}
\sloppy

% Tắt hoàn toàn việc tự chèn gạch nối khi ngắt dòng: tiếng Việt đơn âm tiết
% không bao giờ được tách giữa từ, và babel/vietnam hay ngắt sai từ tiếng Anh.
% Khoảng cách dư do không ngắt được \sloppy + \emergencystretch hấp thụ.
\hyphenpenalty=10000
\exhyphenpenalty=10000

% Cấm ngắt giữa từ tiếng Anh phổ biến (mặc định LaTeX ngắt sai vị trí với
% babel tiếng Việt, vd: "Repos-itory" thay vì "Re-pos-i-to-ry").
\hyphenation{
  Repository Repositories
  Middleware Middlewares
  Controller Controllers
  Service Services
  Prisma Express Helmet
  Frontend Backend
  Database TypeScript JavaScript
  Socket Sequelize
  PostgreSQL MongoDB
  Schema
  Strategy Strategies
  Notification Notifications
  Rotation Rotations
  BKBQP BKTTCP CSTĐCS CSTĐTQ ĐVQT
  HCCSVV HCBVTQ HCQKQT NCKH
  KNCVSNXDQDNDVN QĐNDVN HVKHQS
  ĐATN PM QLKT
  personnelId quanNhan danhHieu
}

\DeclareUnicodeCharacter{2013}{--}
\DeclareUnicodeCharacter{2014}{---}
\DeclareUnicodeCharacter{2018}{`}
\DeclareUnicodeCharacter{2019}{'}
\DeclareUnicodeCharacter{201C}{``}
\DeclareUnicodeCharacter{201D}{''}
\DeclareUnicodeCharacter{2026}{\ldots}
\DeclareUnicodeCharacter{2192}{\ensuremath{\rightarrow}}
\DeclareUnicodeCharacter{2194}{\ensuremath{\leftrightarrow}}
\DeclareUnicodeCharacter{2190}{\ensuremath{\leftarrow}}

% ===================================================


\fancypagestyle{plain}{%
\fancyhf{} % clear all header and footer fields
\fancyfoot[C]{\thepage} % page number centered at bottom
\renewcommand{\headrulewidth}{0pt}
\renewcommand{\footrulewidth}{0pt}}

\setlength{\headheight}{18pt}

\def \TITLE{Xây dựng phần mềm quản lý khen thưởng tại Học viện Khoa học Quân sự}
\def \AUTHOR{Trần Anh Đức}

% ===================================================
\titleformat
    {\chapter} % command
    [hang] % shape
    {\centering\bfseries} % format
    {CHƯƠNG \thechapter.\ } % label
    {0pt} %sep
    {} % before
    [] % after
\titlespacing*{\chapter}{0pt}{-20pt}{20pt}

\titleformat
    {\section} % command
    [hang] % shape
    {\bfseries} % format
    {\thechapter.\arabic{section}\ \ \ \ } % label
    {0pt} %sep
    {} % before
    [] % after
\titlespacing{\section}{0pt}{\parskip}{0.5\parskip}

\titleformat
    {\subsection} % command
    [hang] % shape
    {\bfseries} % format
    {\thechapter.\arabic{section}.\arabic{subsection}\ \ \ \ } % label
    {0pt} %sep
    {} % before
    [] % after
\titlespacing{\subsection}{30pt}{\parskip}{0.5\parskip}

\renewcommand\thesubsubsection{\alph{subsubsection}}
\titleformat
    {\subsubsection} % command
    [hang] % shape
    {\bfseries} % format
    {\alph{subsubsection}, \ } % label
    {0pt} %sep
    {} % before
    [] % after
\titlespacing{\subsubsection}{50pt}{\parskip}{0.5\parskip}

% \newcommand{\titlesize}{\fontsize{18pt}{23pt}\selectfont}
% \newcommand{\subtitlesize}{\fontsize{16pt}{21pt}\selectfont}
% \titleclass{\part}{top}
% \titleformat{\part}[display]
%   {\normalfont\huge\bfseries}{\centering}{20pt}{\Huge\centering}
% \titlespacing{\part}{0pt}{em}{1em}
% \titlespacing{\section}{0pt}{\parskip}{0.5\parskip}
% \titlespacing{\subsection}{0pt}{\parskip}{0.5\parskip}
% \titlespacing{\subsubsection}{0pt}{\parskip}{0.5\parskip}



% ===================================================
\usepackage{hyperref}
\hypersetup{pdfborder = {0 0 0}} %
\hypersetup{pdftitle={\TITLE},
	pdfauthor={\AUTHOR}}
	
\usepackage[all]{hypcap} % Cho phép tham chiếu chính xác đến hình ảnh và bảng biểu

\graphicspath{{Hinhve/}{../Hinhve/}} % Thư mục chứa các hình ảnh

\counterwithin{figure}{chapter} % Đánh số hình ảnh kèm theo chapter. Ví dụ: Hình 1.1, 1.2,..

\title{\bf \TITLE}
\author{\AUTHOR}

\setcounter{secnumdepth}{3} % Cho phép subsubsection trong report
% \setcounter{tocdepth}{3} % Chèn subsubsection vào bảng mục lục

\theoremstyle{definition}
\newtheorem{example}{Ví dụ}[chapter] % Định nghĩa môi trường ví dụ

\onehalfspacing
%Khoảng cách xuống dòng
\setlength{\parskip}{6pt}
%Lùi đầu dòng
\setlength{\parindent}{15pt}



% =========================== BODY ===============
\begin{document}
% \newgeometry{top=2cm, bottom=2cm, left=2cm, right=2cm}
\subfile{Bia} % Phần bìa
\subfile{Bia_lot} % Phần bìa lót
% \restoregeometry

% ===================================================
\pagenumbering{gobble} % Front matter không đánh số trang
\pagestyle{empty} % Header/footer rỗng cho toàn bộ front matter

\newpage
\subfile{Chuong/0_2_Loi_cam_on.tex}

\newpage
\subfile{Chuong/0_3_Tom_tat_noi_dung.tex}

% Tóm tắt tiếng Anh (Abstract) là phần tuỳ chọn theo mẫu SOICT — tạm không trình bày.
%\newpage
%\subfile{Chuong/0_4_Tom_tat_noi_dung_English.tex}

% ===================================================
\renewcommand*\contentsname{MỤC LỤC}

\titlecontents{chapter}
    [0.0cm]             % left margin
    {\bfseries\vspace{0.3cm}}                  % above code
    {{\bfseries{\scshape}
    CHƯƠNG \thecontentslabel.\ }}
    % numbered format
    {}         % unnumbered format
    {\titlerule*[0.3pc]{.}\contentspage}         % filler-page-format, e.g dots

    
\titlecontents{section}
    [0.0cm]             % left margin
    {\vspace{0.3cm}}                  % above code
    {\thecontentslabel \ } % numbered format
    {}         % unnumbered format
    {\titlerule*[0.3pc]{.}\contentspage}         % filler-page-format, e.g dots
    
\titlecontents{subsection}
    [1.0cm]             % left margin
    {\vspace{0.3cm}}                  % above code
    {\thecontentslabel \ } % numbered format
    {}         % unnumbered format
    {\titlerule*[0.3pc]{.}\contentspage}         % filler-page-format, e.g dots

 % Tạo mục lục tự động
\addtocontents{toc}{\protect\thispagestyle{empty}}
\tableofcontents 
\thispagestyle{empty}
\cleardoublepage

% \pagenumbering{roman}
%Tạo danh mục hình vẽ.
\renewcommand{\listfigurename}{DANH MỤC HÌNH VẼ}
{\let\oldnumberline\numberline
\renewcommand{\numberline}{Hình~\oldnumberline}
\listoffigures} 
% \phantomsection\addcontentsline{toc}{section}{\numberline {} DANH MỤC HÌNH VẼ}
\newpage


 %Tạo danh mục bảng biểu.
\renewcommand{\listtablename}{DANH MỤC BẢNG BIỂU}
{\let\oldnumberline\numberline
\renewcommand{\numberline}{Bảng~\oldnumberline}
\listoftables}
% \phantomsection\addcontentsline{toc}{section}{\numberline {} DANH MỤC BẢNG BIỂU}

\clearpage
\subfile{Chuong/0_5_Danh_muc_viet_tat.tex}

\clearpage
\pagenumbering{arabic}

\pagestyle{fancy}
\fancyhf{}
\fancyhead[RE, LO]{\small\leftmark}
%\fancyhead[LE]{\rightmark}
\fancyfoot[RE, LO]{\thepage}

\chapter{GIỚI THIỆU ĐỀ TÀI}
\label{chapter:Introduction}
\subfile{Chuong/1_Gioi_thieu} % Phần mở đầu

\newpage
%\pagestyle{fancy} % Áp dụng header và footer
\chapter{KHẢO SÁT VÀ PHÂN TÍCH YÊU CẦU}
\label{chapter:Related_works}
\subfile{Chuong/2_Khao_sat}


\newpage
%\pagestyle{fancy} % Áp dụng header và footer
\chapter{NỀN TẢNG LÝ THUYẾT VÀ CÔNG NGHỆ SỬ DỤNG}
\label{chapter:Methodology}
\subfile{Chuong/3_Cong_nghe}

\newpage
%\pagestyle{fancy} % Áp dụng header và footer
\chapter{PHÂN TÍCH THIẾT KẾ, TRIỂN KHAI VÀ ĐÁNH GIÁ HỆ~THỐNG}
\label{chapter:Experiment}
\subfile{Chuong/4_Ket_qua_thuc_nghiem}

\newpage
%\pagestyle{fancy} % Áp dụng header và footer
\chapter{CÁC GIẢI PHÁP VÀ ĐÓNG GÓP NỔI BẬT}
\label{chapter:SolutionAndContribution}
\subfile{Chuong/5_Giai_phap_dong_gop}
\newpage
%\pagestyle{fancy} % Áp dụng header và footer
\chapter{KẾT LUẬN VÀ HƯỚNG PHÁT TRIỂN}
\label{chapter:conclusion}
\subfile{Chuong/6_Ket_luan}

% ===================================================
\newpage
\renewcommand\bibname{TÀI LIỆU THAM KHẢO}
\begin{thebibliography}{99}
\phantomsection\addcontentsline{toc}{chapter}{TÀI LIỆU THAM KHẢO}
\bibitem{luat062022} Quốc hội nước CHXHCN Việt Nam, \textit{Luật Thi đua, Khen thưởng số 06/2022/QH15}, 2022.
\bibitem{nextjs} Next.js. [Online]. Available: \url{https://nextjs.org} (visited on 04/03/2026).
\bibitem{react} React. [Online]. Available: \url{https://react.dev} (visited on 04/03/2026).
\bibitem{express} Express. [Online]. Available: \url{https://expressjs.com} (visited on 05/03/2026).
\bibitem{nodejs} Node.js. [Online]. Available: \url{https://nodejs.org} (visited on 05/03/2026).
\bibitem{typescript} TypeScript. [Online]. Available: \url{https://www.typescriptlang.org} (visited on 06/03/2026).
\bibitem{postgresql} PostgreSQL. [Online]. Available: \url{https://www.postgresql.org} (visited on 08/03/2026).
\bibitem{prisma} Prisma. [Online]. Available: \url{https://www.prisma.io} (visited on 10/03/2026).
\bibitem{antdesign} Ant Design. [Online]. Available: \url{https://ant.design} (visited on 12/03/2026).
\bibitem{tailwind} Tailwind CSS. [Online]. Available: \url{https://tailwindcss.com} (visited on 13/03/2026).
\bibitem{socketio} Socket.IO. [Online]. Available: \url{https://socket.io} (visited on 16/03/2026).
\bibitem{jwt} JSON Web Token (JWT). [Online]. Available: \url{https://jwt.io} (visited on 18/03/2026).
\bibitem{zod} Zod. [Online]. Available: \url{https://zod.dev} (visited on 20/03/2026).
\bibitem{exceljs} ExcelJS. [Online]. Available: \url{https://github.com/exceljs/exceljs} (visited on 22/03/2026).
\bibitem{layeredref} Microsoft, N-tier architecture style. [Online]. Available: \url{https://learn.microsoft.com/azure/architecture/guide/architecture-styles/n-tier} (visited on 02/03/2026).
\end{thebibliography}

\end{document}
